% arXiv can compile this with pdfLaTeX; add \pdfoutput=1 if an arXiv
% upload specifically requests it for bitmap figures. It is omitted here so
% the bundled Tectonic/XeTeX compiler can build the local preview.
\documentclass[11pt]{article}

\usepackage[margin=1in]{geometry}
\usepackage{cmap}
\usepackage[T1]{fontenc}
\usepackage[utf8]{inputenc}
\usepackage{microtype}
\usepackage{amsmath,amssymb,amsthm}
\usepackage{booktabs}
\usepackage{graphicx}
\usepackage{enumitem}
\usepackage[numbers,sort&compress]{natbib}
\usepackage[colorlinks=true,allcolors=blue]{hyperref}

\newtheorem{proposition}{Proposition}
\newtheorem{assumption}{Assumption}

\title{Foveated Local Loops with HydraLoRA:\\
A VRAM-Constrained Vision World Model for Local Evidence Acquisition}

\author{Author information omitted for draft}
\date{\today}

\begin{document}
\maketitle

\begin{abstract}
Vision world models are usually discussed as compact latent predictors, yet real visual reasoning often depends on small local evidence: a pin label, a trace junction, a polarity marker, or a subtle defect.
Processing every pixel at uniformly high resolution is wasteful under consumer-GPU memory constraints, while processing only crops discards the global context that makes local evidence meaningful.
We propose a foveated world-model architecture that separates a low-resolution global peripheral loop from high-resolution local loops.
The global loop maintains a predictive latent state, coarse topology, uncertainty, and task context; the local loops inspect a small set of regions of interest (ROIs), extract typed evidence, and update a graph-like memory.
To make local visual specialization practical, we reorganize a Simula-trained Local LoRA Bank into a HydraLoRA-style shared-$A$/multi-$B$ expert structure with certified bundle routing.
We provide a mathematical formulation for token-aware VRAM budgeting, adapter parameter savings, top-$k$ expert routing, residual-driven synthetic curriculum construction, and graph-memory updates.
The paper is positioned as a compact research proposal and experimental protocol for circuit-diagram and PCB understanding, a domain where global topology and tiny local cues interact strongly.
\end{abstract}

\section{Introduction}

The motivating question is simple: can a vision world model keep the whole scene in memory while spending high-resolution computation only where evidence is needed?
This question matters in domains such as circuit diagrams and PCB inspection, where a low-resolution image may preserve the global layout but erase the small evidence that determines graph structure, polarity, part identity, or defect status.
A uniform high-resolution pass is expensive in visual tokens, activation memory, and key-value (KV) cache.
A crop-only pipeline is cheaper but risks losing the peripheral context required to decide which crop matters and how its evidence should update the world state.

We propose \emph{Foveated Local Loops with HydraLoRA}, a VRAM-aware perception architecture.
The system has a low-resolution global loop, a high-recall ROI proposal loop, a high-resolution local loop, a HydraLoRA-based local-skill router, a typed evidence extractor, and a graph-memory update loop.
The architecture draws conceptual support from world models and JEPA-style latent prediction \citep{ha2018worldmodels,hafner2023dreamerv3,assran2023ijepa,bardes2024vjepa,maes2026leworldmodel}, low-rank adaptation \citep{hu2021lora,dettmers2023qlora}, HydraLoRA \citep{tian2024hydralora}, and reasoning-driven synthetic data generation \citep{davidson2026simula}.

The main contribution of this draft is not a new empirical result.
Instead, it turns the existing architecture notes into a mathematically specified workshop-style paper:
\begin{enumerate}[leftmargin=*]
  \item a formal foveated observation model for global state, ROI crops, typed evidence, and graph memory;
  \item a token-aware VRAM budget that separates adapter memory from visual-token and KV-cache costs;
  \item a HydraLoRA memory analysis for independent LoRA banks, shared-$A$ banks, and active top-$k$ routing;
  \item a residual-driven objective that maps world-model prediction failures into Simula-style local training data;
  \item an experimental protocol for circuit/PCB understanding under fixed memory budgets.
\end{enumerate}

\begin{figure}[t]
  \centering
  \setlength{\fboxsep}{8pt}
  \fbox{%
    \begin{minipage}{0.92\linewidth}
      \centering
      Low-res global loop
      $\rightarrow$
      ROI proposal
      $\rightarrow$
      high-res local crop
      $\rightarrow$
      HydraLoRA skill routing
      $\rightarrow$
      typed evidence
      $\rightarrow$
      graph/world-state update
      $\rightarrow$
      feedback to ROI proposal
      \vspace{0.35em}

      \small
      Offline Simula skill forge:
      residual mining
      $\rightarrow$
      taxonomy update
      $\rightarrow$
      synthetic ROI tasks
      $\rightarrow$
      critic filtering
      $\rightarrow$
      adapter training
      $\rightarrow$
      certified bundle registry.
    \end{minipage}
  }
  \caption{High-level foveated loop used by the draft system: a low-resolution global state proposes local inspections, high-resolution crops produce evidence, and the evidence returns to memory and future ROI selection.}
  \label{fig:foveated-loop}
\end{figure}

\section{Related Work}

\paragraph{World models and latent prediction.}
World models learn compact internal dynamics for perception and planning \citep{ha2018worldmodels,hafner2023dreamerv3}.
JEPA-style methods predict in representation space rather than pixel space \citep{assran2023ijepa,bardes2024vjepa}.
LeWorldModel (LeWM) is especially relevant because it trains a JEPA-style latent world model from pixels with a compact prediction objective and a Gaussian latent regularizer \citep{maes2026leworldmodel}.
Our global loop borrows this latent-prediction view, but the target task here is local evidence acquisition in static or semi-static visual artifacts, not direct action-conditioned control.

\paragraph{LoRA banks and HydraLoRA.}
LoRA freezes a base weight matrix $W_0$ and learns a low-rank update $\Delta W = BA$ \citep{hu2021lora}.
QLoRA further shows that quantized base weights and low-rank adapters can make training feasible under tight GPU memory budgets \citep{dettmers2023qlora}.
However, a multi-skill LoRA bank introduces a new problem: even if each adapter is small, many adapters create resident-memory, loading, routing, and compatibility costs.
HydraLoRA proposes an asymmetric shared-$A$/multi-$B$ structure for efficient fine-tuning \citep{tian2024hydralora}.
We reinterpret that structure as an organizer for local visual skills.

\paragraph{Synthetic curricula and local skill construction.}
Simula frames synthetic data generation as a controllable mechanism design problem with taxonomy construction, global and local diversification, complexification, and critic filtering \citep{davidson2026simula}.
In our setting, Simula is not a runtime reasoner.
It is an offline skill forge that converts failure logs, graph conflicts, and uncertain ROI cases into local training tuples for adapters and routers.

\section{Problem Formulation}

Let $I_t \in \mathbb{R}^{H \times W \times 3}$ be an image or video frame at step $t$.
For static circuit/PCB analysis, $t$ can index inspection iterations rather than physical time.
The system keeps a global latent state $z_t \in \mathbb{R}^{d_z}$, a graph memory $G_t=(V_t,E_t,\psi_t)$, and an evidence table $\mathcal{E}_t$.
Here $\psi_t$ stores beliefs over graph attributes such as component labels, pin connections, junction states, polarity, and defect annotations.

\paragraph{Foveated observation.}
The global loop observes a low-resolution image
\begin{equation}
  x_t^g = D_\rho(I_t),
\end{equation}
where $D_\rho$ downsamples by factor $\rho$.
The ROI loop selects at most $K$ boxes
\begin{equation}
  \mathcal{R}_t = \{r_{t,1},\ldots,r_{t,k_t}\}, \qquad k_t \le K,
\end{equation}
and extracts high-resolution crops
\begin{equation}
  x_{t,j}^{\ell} = C(I_t,r_{t,j},m_j),
\end{equation}
where $m_j$ is an optional context margin.
The full foveated observation is
\begin{equation}
  o_t = \left(x_t^g,\{x_{t,j}^{\ell}\}_{j=1}^{k_t},G_t,\mathcal{E}_t\right).
\end{equation}

\paragraph{Global latent prediction.}
The global encoder $E_g$ and predictor $P$ define
\begin{equation}
  z_t = E_g(x_t^g,G_t), \qquad
  \hat z_{t+1}=P(z_t,u_t),
\end{equation}
where $u_t$ is the current query, inspection action, or task context.
Following the LeWM-style intuition, the global training objective includes a prediction term and a distributional anti-collapse regularizer:
\begin{equation}
  \mathcal{L}_{\mathrm{global}}
  =
  \left\|\hat z_{t+1} - \mathrm{sg}(z_{t+1})\right\|_2^2
  +
  \lambda_{\mathrm{sig}}\mathcal{D}_{\mathrm{sig}}(\{z_t\},\mathcal{N}(0,I)).
  \label{eq:global-loss}
\end{equation}
The operator $\mathrm{sg}$ denotes stop-gradient when used in a teacher-target form.
The regularizer $\mathcal{D}_{\mathrm{sig}}$ can be instantiated by random one-dimensional projections, comparing projected latent samples to a standard normal distribution.
For this proposal, the exact LeWM regularizer is not re-derived; the important modeling role is to prevent collapsed global states while preserving a compact predictive latent.

\paragraph{Residuals as local data signals.}
A prediction residual is not only a scalar loss.
We define a structured residual
\begin{equation}
  \delta_t =
  \left[
  d_z(\hat z_{t+1},z_{t+1}),
  d_G(\hat G_{t+1},G_{t+1}),
  1-\bar c_t,
  \chi_{\mathrm{conflict}}(G_t,\mathcal{E}_t)
  \right],
  \label{eq:residual}
\end{equation}
where $d_z$ measures latent mismatch, $d_G$ measures graph inconsistency, $\bar c_t$ is mean evidence confidence, and $\chi_{\mathrm{conflict}}$ flags incompatible local claims.
The offline Simula loop maps $\delta_t$ to taxonomy nodes, hard negatives, ROI labels, and required local skills.

\section{Method}

\subsection{Budget-Aware ROI Selection}

Let $T_g$ be the number of visual tokens for the low-resolution global image and $T_j$ the token count for ROI crop $j$.
For patch size $p$ and crop dimensions $h_j \times w_j$,
\begin{equation}
  T_j \approx \left\lceil \frac{h_j}{p} \right\rceil
          \left\lceil \frac{w_j}{p} \right\rceil.
\end{equation}
The total visual-token count is
\begin{equation}
  T_{\mathrm{vis}} = T_g + \sum_{j=1}^{k_t} T_j.
  \label{eq:tokens}
\end{equation}

The ROI scheduler assigns utility
\begin{equation}
  U(r)
  =
  \alpha u_{\mathrm{unc}}(r)
  + \beta u_{\mathrm{task}}(r)
  + \gamma u_{\mathrm{graph}}(r)
  + \xi u_{\mathrm{failure}}(r)
  - \eta c_{\mathrm{tok}}(r)
  - \mu c_{\mathrm{adapt}}(r),
  \label{eq:roi-utility}
\end{equation}
where the positive terms represent uncertainty, task relevance, graph conflict, and historical failure signals, while the negative terms estimate token and adapter costs.
The selected ROI set solves a small budgeted selection problem:
\begin{equation}
  \mathcal{R}_t^\star
  =
  \arg\max_{\mathcal{R}\subseteq \mathcal{C}_t}
  \sum_{r\in\mathcal{R}} U(r)
  \quad
  \mathrm{s.t.}
  \quad
  |\mathcal{R}|\le K,\quad
  T_g+\sum_{r\in\mathcal{R}}T(r)\le T_{\max}.
  \label{eq:roi-select}
\end{equation}
Here $\mathcal{C}_t$ is a high-recall candidate set.
In the MVP design, $K=3$ follows the local orchestration notes.

\subsection{VRAM Accounting}

Peak memory is approximated as
\begin{equation}
  M_{\mathrm{peak}}
  \approx
  M_{\mathrm{base}}(q)
  + M_{\mathrm{adapt}}(\mathcal{B}_t)
  + M_{\mathrm{KV}}(T_{\mathrm{vis}})
  + M_{\mathrm{act}}(B,T_{\mathrm{vis}})
  + M_{\mathrm{misc}},
  \label{eq:vram}
\end{equation}
where $q$ is base-model precision, $\mathcal{B}_t$ is the active certified adapter bundle, $B$ is batch size, and $M_{\mathrm{misc}}$ includes runtime allocator overhead.
For a transformer with $L$ layers, hidden dimension $d$, batch size $B$, and KV precision $b_{\mathrm{kv}}$ bytes, a simple KV-cache estimate is
\begin{equation}
  M_{\mathrm{KV}}(T_{\mathrm{vis}})
  \approx
  2 L B T_{\mathrm{vis}} d\, b_{\mathrm{kv}}.
  \label{eq:kv}
\end{equation}
This term is why the proposal treats visual-token routing as first-class: reducing adapter weights alone may not reduce peak VRAM if $T_{\mathrm{vis}}$ is large.

\begin{proposition}[Foveated token ratio]
For an image processed at full resolution $H\times W$ and patch size $p$, the approximate token count is $T_{\mathrm{full}}\approx HW/p^2$.
If the foveated system uses a global image of size $H_g\times W_g$ and $K$ ROI crops of size $h\times w$, then
\begin{equation}
  \frac{T_{\mathrm{fov}}}{T_{\mathrm{full}}}
  \approx
  \frac{H_gW_g + Khw}{HW}.
  \label{eq:fov-ratio}
\end{equation}
\end{proposition}

\begin{proof}
Using the same patch size $p$, the global token count is approximately $H_gW_g/p^2$ and each crop contributes $hw/p^2$ tokens.
Dividing $T_{\mathrm{fov}}\approx (H_gW_g+Khw)/p^2$ by $T_{\mathrm{full}}\approx HW/p^2$ gives Eq.~\ref{eq:fov-ratio}.
\end{proof}

\subsection{HydraLoRA Local Skill Bank}

For a frozen linear layer $W_0\in\mathbb{R}^{d_{\mathrm{out}}\times d_{\mathrm{in}}}$, an independent LoRA expert $i$ uses
\begin{equation}
  W_i = W_0 + \frac{\alpha}{r}B_iA_i,
  \qquad
  A_i\in\mathbb{R}^{r\times d_{\mathrm{in}}},
  \quad
  B_i\in\mathbb{R}^{d_{\mathrm{out}}\times r}.
\end{equation}
HydraLoRA shares $A$ and keeps skill-specific $B_i$ matrices:
\begin{equation}
  W(x)=W_0+\frac{\alpha}{r}\sum_{i\in S_k(x)}g_i(x)B_iA,
  \label{eq:hydra}
\end{equation}
where $g_i(x)$ is a router score and $S_k(x)$ is the top-$k$ expert set.
The router input combines global context, ROI embedding, ROI type hint, task instruction, complexity estimate, and registry metadata:
\begin{equation}
  g(x)=\mathrm{softmax}\left(R_\phi([z_t,h_j,\tau_j,q_t,c_j])\right).
\end{equation}

\begin{proposition}[Adapter parameter savings]
For one adapted layer with $N$ experts and rank $r$, an independent LoRA bank has
\begin{equation}
  P_{\mathrm{bank}} = N r(d_{\mathrm{in}}+d_{\mathrm{out}})
\end{equation}
parameters, while a shared-$A$ HydraLoRA bank has
\begin{equation}
  P_{\mathrm{hydra}} = r d_{\mathrm{in}} + N r d_{\mathrm{out}}.
\end{equation}
If $d_{\mathrm{in}}=d_{\mathrm{out}}=d$, then
\begin{equation}
  \frac{P_{\mathrm{hydra}}}{P_{\mathrm{bank}}}=\frac{N+1}{2N}.
  \label{eq:hydra-ratio}
\end{equation}
If only $k$ of the $N$ $B$ experts are active on GPU, the active adapter parameter ratio relative to a fully resident independent bank is
\begin{equation}
  \frac{P_{\mathrm{active}}}{P_{\mathrm{bank}}}=\frac{k+1}{2N}.
  \label{eq:active-ratio}
\end{equation}
\end{proposition}

\begin{proof}
The independent bank stores $N$ copies of $A_i$ and $B_i$, giving $Nr d_{\mathrm{in}}+Nr d_{\mathrm{out}}$.
The Hydra bank stores one $A$ and $N$ $B_i$ matrices, giving $rd_{\mathrm{in}}+Nr d_{\mathrm{out}}$.
The square-layer ratio follows by substituting $d_{\mathrm{in}}=d_{\mathrm{out}}=d$.
For active top-$k$ loading, one shared $A$ and $k$ $B_i$ matrices are resident, giving $rd+krd$ parameters and Eq.~\ref{eq:active-ratio}.
\end{proof}

\subsection{Router and Adapter Training}

Each synthetic or real training item is represented as
\begin{equation}
  s =
  (x^g,x^\ell,r,\tau,y,e,\Delta G,m),
\end{equation}
where $\tau$ is a Simula taxonomy path, $y$ is a task label, $e$ is typed local evidence, $\Delta G$ is the graph update label, and $m$ contains quality, complexity, and critic metadata.
The total training objective is
\begin{align}
  \mathcal{L}
  =&\;
  \mathcal{L}_{\mathrm{evidence}}
  + \lambda_G \mathcal{L}_{\mathrm{graph}}
  + \lambda_R \mathcal{L}_{\mathrm{router}}
  + \lambda_B \mathcal{L}_{\mathrm{balance}}
  + \lambda_H \mathcal{L}_{\mathrm{entropy}}
  \nonumber\\
  &+
  \lambda_D \mathcal{L}_{\mathrm{distill}}
  + \lambda_M \left[\max(0,M_{\mathrm{peak}}-M_{\mathrm{budget}})\right]^2.
  \label{eq:total-loss}
\end{align}

The router can be supervised by taxonomy-derived expert labels $q_i(s)$:
\begin{equation}
  \mathcal{L}_{\mathrm{router}}
  =
  -\sum_i q_i(s)\log g_i(s).
\end{equation}
To reduce expert collapse, we use a load-balancing term
\begin{equation}
  \mathcal{L}_{\mathrm{balance}}
  =
  N\sum_{i=1}^{N}
  \left(\frac{1}{|\mathcal{D}|}\sum_{s\in\mathcal{D}}g_i(s)-\frac{1}{N}\right)^2.
\end{equation}
To let difficult samples use more experts while keeping easy samples cheap, we set a target entropy from normalized complexity $c(s)\in[0,1]$:
\begin{equation}
  H^\star(s)=H_{\min}+(H_{\max}-H_{\min})c(s),
  \qquad
  \mathcal{L}_{\mathrm{entropy}}=(H(g(s))-H^\star(s))^2.
\end{equation}
If a teacher independent LoRA bank exists, shared-$A$ distillation can be added:
\begin{equation}
  \mathcal{L}_{\mathrm{distill}}
  =
  \left\|
  f_{\mathrm{hydra}}(x)-f_{\mathrm{bank}}(x)
  \right\|_2^2.
\end{equation}

\subsection{Graph Memory Update}

Local outputs are not treated as free-form text.
They become typed evidence items
\begin{equation}
  e_j=(a_j,v_j,b_j,c_j,\mathcal{B}_j),
\end{equation}
where $a_j$ is an attribute or graph claim, $v_j$ its value, $b_j$ the bounding region, $c_j$ confidence, and $\mathcal{B}_j$ the adapter bundle used.
For a graph proposition $y$ such as ``pin 1 connects to net N'', a log-odds update is
\begin{equation}
  \ell_{t+1}(y)
  =
  \ell_t(y)
  +
  \sum_{e_j\in\mathcal{E}_{t+1}}
  w(e_j,y)
  \log
  \frac{p(e_j\mid y)}{p(e_j\mid \neg y)}
  -
  \Omega_{\mathrm{conflict}}(y,G_t).
  \label{eq:graph-update}
\end{equation}
The belief is $b_{t+1}(y)=\sigma(\ell_{t+1}(y))$.
Conflicts do not immediately force a conclusion; they raise future ROI utility through $u_{\mathrm{graph}}(r)$ in Eq.~\ref{eq:roi-utility}.

\begin{proposition}[Calibrated evidence increases belief]
If $\Omega_{\mathrm{conflict}}=0$, $w(e,y)>0$, and $p(e\mid y)>p(e\mid\neg y)$, then observing $e$ strictly increases the log-odds and belief of $y$.
\end{proposition}

\begin{proof}
Under the stated conditions, the log-likelihood ratio in Eq.~\ref{eq:graph-update} is positive.
Therefore $\ell_{t+1}(y)>\ell_t(y)$.
Because the sigmoid function is strictly increasing, $b_{t+1}(y)>b_t(y)$.
\end{proof}

\subsection{Certified Bundle Registry}

The runtime system is not allowed to combine arbitrary adapters or sweep scales.
Let $b$ be a candidate bundle and let $\mathcal{V}$ be a validation set.
We register $b$ only if
\begin{align}
  \mathrm{Score}(b;\mathcal{V}) &\ge \mathrm{Score}(\mathrm{base};\mathcal{V})+\epsilon_{\mathrm{task}},\\
  \mathrm{Regress}(b;\mathcal{V}_{\mathrm{guard}}) &\le \epsilon_{\mathrm{reg}},\\
  M_{\mathrm{peak}}(b) &\le M_{\mathrm{budget}},\\
  L_{95}(b) &\le L_{\mathrm{budget}},\\
  \mathrm{ConflictRate}(b) &\le \epsilon_{\mathrm{conflict}}.
\end{align}
Bundles failing these gates are quarantined rather than silently deleted.
At inference time, router abstention is allowed:
\begin{equation}
  \max_i g_i(x)<\theta_{\mathrm{route}}
  \quad\Rightarrow\quad
  \mathrm{use\;base/generic\;fallback/human\;review}.
\end{equation}

\section{Offline Simula Skill Forge}

The offline loop converts residuals into training data.
Let $\tau$ be a taxonomy node and $D_t(\tau)$ the current training distribution over local cases.
We update sampling priority from failure residuals:
\begin{equation}
  \pi_{t+1}(\tau)
  =
  \frac{
  \pi_t(\tau)\exp(\rho\,\bar\delta_t(\tau))
  }{
  \sum_{\tau'}\pi_t(\tau')\exp(\rho\,\bar\delta_t(\tau'))
  },
  \label{eq:curriculum}
\end{equation}
where $\bar\delta_t(\tau)$ aggregates residual severity for examples mapped to node $\tau$.
Synthetic data generation then targets high-priority nodes with local diversification, complexity control, hard negatives, and double-critic filtering.
The resulting data includes ROI labels as first-class supervision, not merely final answers.

\section{Experimental Protocol}

\paragraph{Domain.}
The primary evaluation domain is circuit diagrams and PCB imagery.
This domain is graph-like, visually dense, and sensitive to tiny local evidence.
Document and chart images can be included as secondary sanity checks because the orchestration notes already define document extraction and chart reading as neighboring skill families.

\paragraph{Baselines.}
We propose the following comparison set:
\begin{enumerate}[leftmargin=*]
  \item base VLM/world model, no local adapters;
  \item full high-resolution input with no foveated routing;
  \item foveated ROI routing without LoRA;
  \item single monolithic LoRA;
  \item independent skill-specific LoRA bank;
  \item HydraLoRA shared-$A$ bank;
  \item HydraLoRA with Simula taxonomy router;
  \item HydraLoRA with visual-budget-aware top-$k$ active $B$ loading.
\end{enumerate}

\paragraph{Metrics.}
Task quality should be measured by accuracy, macro F1, OCR accuracy, per-taxonomy accuracy, graph-edge F1, conflict resolution rate, and OOD accuracy.
Memory and serving quality should be measured by adapter parameters, adapter resident VRAM, visual-token count, peak allocated VRAM, average VRAM, latency p50/p95, adapter load time, and router overhead.
Router quality should be measured by router accuracy, top-$k$ hit rate, gate entropy, expert utilization, collapse rate, abstention rate, and fallback success rate.

\paragraph{Ablations.}
The key ablations are $N\in\{2,4,8,16\}$ experts, LoRA rank $r\in\{4,8,16,32\}$, top-$k\in\{1,2,4,\mathrm{all}\}$, with/without Simula complexity metadata, with/without load balancing, with/without certified registry rehearsal, and full-image high-resolution input versus ROI-routed input.

\begin{table}[t]
  \centering
  \caption{Idealized square-layer adapter ratios. These ratios exclude activations, visual tokens, KV cache, and runtime allocator overhead.}
  \label{tab:ratios}
  \begin{tabular}{cccc}
    \toprule
    Experts $N$ & Full Hydra ratio $(N+1)/(2N)$ & Active top-1 ratio & Active top-2 ratio \\
    \midrule
    2  & 0.7500 & 0.5000 & 0.7500 \\
    4  & 0.6250 & 0.2500 & 0.3750 \\
    8  & 0.5625 & 0.1250 & 0.1875 \\
    16 & 0.5313 & 0.0625 & 0.0938 \\
    \bottomrule
  \end{tabular}
\end{table}

\section{Discussion}

The mathematical analysis clarifies two claims that were implicit in the design notes.
First, HydraLoRA mainly reduces adapter-bank memory and active adapter memory; it does not by itself guarantee low peak VRAM.
Peak VRAM also depends on high-resolution visual tokens, activations, and KV cache.
Therefore, the foveated token policy and the adapter policy must be evaluated separately.
Second, the local loop is not a crop classifier.
Its output is typed evidence that changes graph beliefs and future ROI utilities.
This is why the architecture is better described as local evidence acquisition for a world model than as OCR plus adapters.

The proposal also has clear risks.
The A/B asymmetry observed in LLM HydraLoRA may be weaker in small vision or vision-language backbones.
Synthetic data can induce shortcuts if taxonomy labels leak into visual style.
Wrong expert routing may create more confident errors than base-only inference.
Finally, if the base model lacks sufficient OCR or grounding ability, local adapters may not recover the missing capability.
The certified registry, abstention, hard negatives, guard validation, and quarantine policy are meant to make these risks measurable rather than hidden.

\section{Conclusion}

We presented a mathematically specified draft for a VRAM-constrained foveated vision world model.
The system preserves a low-resolution predictive global state, inspects a few high-resolution ROIs, routes local skill adapters through a HydraLoRA-style shared-$A$/multi-$B$ bank, extracts typed evidence, and updates graph memory.
The strengthened formulation adds token-aware VRAM bounds, adapter memory propositions, router objectives, graph-belief updates, and Simula-driven residual curricula.
This gives the project a clearer arXiv-style shape: a compact position paper and experimental protocol for local evidence acquisition under memory constraints.

\bibliographystyle{unsrtnat}
\bibliography{references}

\end{document}
